\section{Introduction}
Hyperspectral imaging (HSI) provides densely sampled spectral observations that support fine-grained land-cover mapping. In contrast to RGB or multispectral imagery, each HSI pixel contains a high-dimensional spectral signature, allowing classification systems to separate materials with subtle reflectance differences. This property makes HSI useful in agriculture, urban mapping, mineral exploration, forestry, water monitoring, and environmental assessment. However, the same high dimensionality that provides discriminative information also introduces redundancy, noise, and severe annotation difficulty.

Deep learning has become the dominant direction for HSI classification. Three-dimensional convolutional neural networks learn joint spectral-spatial features, hybrid CNNs reduce computational cost by separating spectral and spatial processing, and residual or densely connected networks improve feature propagation. Transformer-based methods further model long-range dependencies by self-attention, while graph neural networks represent relational structure among pixels, superpixels, or regions. More recently, state-space models related to Mamba have become attractive because they can process long spectral sequences with linear complexity. These advances have produced strong results, but several practical issues remain insufficiently studied.

First, public HSI datasets usually contain very few labeled pixels compared with the number of spectral bands and spatial samples. Second, class distributions are often imbalanced, so high overall accuracy can hide weak performance on rare land-cover categories. Third, class boundaries and mixed pixels create ambiguous samples whose labels are intrinsically difficult to predict. Fourth, reported performance can vary strongly across random training splits. A method that obtains high accuracy on one split may be less reliable when the available labeled samples change. Therefore, a journal-ready HSI classification study should evaluate not only average accuracy but also stability, rare-class behavior, uncertainty, and reproducibility.

This work studies hyperspectral image classification from the perspective of reliability under label scarcity. The central question is whether graph-state-space modeling with explicit robust optimization can improve rare-class behavior, worst-class performance, and split stability when only a small number of labeled pixels are available. We propose DR-GSMamba, a graph-state-space framework that couples efficient spectral sequence modeling with local graph reasoning and a robust learning objective. The method receives a center spectrum and a local patch for each labeled pixel. A state-space spectral encoder captures long-range band relationships, a compact spatial convolutional stem extracts local texture, and a patch graph branch aggregates non-local contextual information. These representations are fused and classified using both a linear classifier and a prototype uncertainty head.

The main contributions are summarized as follows.
\begin{itemize}
\item We propose a reliability-oriented graph-state-space architecture for label-scarce HSI classification, combining spectral sequence modeling, spatial convolution, and patch-level graph message passing.
\item We introduce an explicit robust training objective that combines sample-level CVaR and class-level worst-risk optimization, making the distributionally robust claim experimentally testable rather than merely descriptive.
\item We incorporate prototype and uncertainty-aware learning to encourage stable class representations and to analyze ambiguous mixed or boundary pixels.
\item We provide a reproducible evaluation protocol with low-label sweeps, multi-seed reporting, rare-class metrics, worst-class metrics, calibration analysis, and ablations that isolate every named component.
\end{itemize}

The remainder of this paper is organized as follows. Section II reviews related work. Section III presents the proposed DR-GSMamba method. Section IV describes the experimental protocol and reproducibility pipeline. Section V concludes the paper and discusses the next benchmark stage.
